\chapter{星积、Moyal 括号与算符对应}
\label{ch:moyal}

\begin{learningobjectives}
完成本章后，读者应能：
\begin{enumerate}
  \item 用星积计算算符乘积的 Weyl 符号；
  \item 从 von Neumann 方程得到 Moyal 演化方程；
  \item 识别经典 Liouville 项和领先量子修正；
  \item 使用 Bopp 算符把主方程转换成相空间偏微分方程。
\end{enumerate}
\end{learningobjectives}

本章的实变量公式继续使用有量纲正则对 $(x,p_x)$，满足 $[\hat x,\hat p_x]=\ii\hbar$；进入复变量公式时，则使用第2章的无量纲 $\alpha=(Q+\ii P)/\sqrt2$。
星积与 Moyal 括号的原始表述见\cite{Moyal1949}；本章固定采用第2章的
归一化，避免把不同文献中的 $\hbar$ 与面积元约定混用。

\section{普通乘积为什么不够}

Weyl 映射保持线性，却不把算符乘积直接映射为普通函数乘积。记
\begin{equation*}
  \hat A\leftrightarrow A_W(x,p_x),
  \qquad
  \hat B\leftrightarrow B_W(x,p_x),
\end{equation*}
则
\begin{equation}
  (\hat A\hat B)_W=A_W\star B_W,
  \label{eq:star-product-correspondence}
\end{equation}
其中星积定义为
\begin{equation}
  A\star B
  =A\exp\!\left[
  \frac{\ii\hbar}{2}
  \left(
  \overleftarrow{\partial_x}\overrightarrow{\partial_{p_x}}
  -\overleftarrow{\partial_{p_x}}\overrightarrow{\partial_x}
  \right)\right]B.
  \label{eq:star-product}
\end{equation}
箭头说明导数作用在哪一侧。定义双向微分算符
\begin{equation}
  \Lambda=
  \overleftarrow{\partial_x}\overrightarrow{\partial_{p_x}}
  -\overleftarrow{\partial_{p_x}}\overrightarrow{\partial_x},
\end{equation}
则 $A\star B=A\exp(\ii\hbar\Lambda/2)B$。

把指数展开，得到
\begin{equation}
  A\star B=AB+\frac{\ii\hbar}{2}\{A,B\}_{\mathrm{PB}}
  +\order(\hbar^2),
  \label{eq:star-expansion}
\end{equation}
其中
\begin{equation}
  \{A,B\}_{\mathrm{PB}}
  =\pder{A}{x}\pder{B}{p_x}
  -\pder{A}{p_x}\pder{B}{x}
\end{equation}
是 Poisson 括号。星积由普通乘积、Poisson 几何和更高阶导数共同组成。

\begin{example}
取 $A=x$、$B=p_x$。因为二阶以上导数全部为零，
\begin{equation}
  x\star p_x=x p_x+\frac{\ii\hbar}{2},
  \qquad
  p_x\star x=p_x x-\frac{\ii\hbar}{2}.
\end{equation}
因此 $(\comm{\hat x}{\hat p_x})_W=x\star p_x-p_x\star x=\ii\hbar$。不对易性没有消失，而是进入星积。
\end{example}

\section{Moyal 括号}

算符对易子的 Weyl 符号定义 Moyal 括号：
\begin{align}
  \{A,B\}_{\mathrm{M}}
  &\equiv\frac{1}{\ii\hbar}(A\star B-B\star A)\\
  &=\frac{2}{\hbar}A
  \sin\!\left(\frac{\hbar}{2}\Lambda\right)B.
  \label{eq:moyal-bracket}
\end{align}
展开正弦可见
\begin{equation}
  \{A,B\}_{\mathrm{M}}
  =\{A,B\}_{\mathrm{PB}}
  -\frac{\hbar^2}{24}A\Lambda^3B
  +\order(\hbar^4).
  \label{eq:moyal-expansion}
\end{equation}
经典极限不是简单令 $\hbar=0$，而是要求高阶导数项相对于 Poisson 项可忽略。若 Wigner 函数在相空间中形成很细的振荡结构，即使 $\hbar$ 数值很小，高阶导数也可能变得重要。

\section{从 von Neumann 方程到 Wigner 演化}

闭合系统的密度算符满足
\begin{equation}
  \pder{\rhohat}{t}=-\frac{\ii}{\hbar}\comm{\Hhat}{\rhohat}.
  \label{eq:von-neumann}
\end{equation}
对\cref{eq:von-neumann} 作 Weyl 变换，利用\cref{eq:star-product-correspondence}，得到精确的 Moyal 方程
\begin{equation}
  \pder{W}{t}=\{H_W,W\}_{\mathrm{M}}.
  \label{eq:moyal-evolution}
\end{equation}
第一项是经典 Liouville 流，后续奇数阶导数包含量子修正。

对
\begin{equation}
  H_W(x,p_x)=\frac{p_x^2}{2m}+V(x),
\end{equation}
\cref{eq:moyal-evolution} 可写成
\begin{equation}
  \pder{W}{t}
  =-\frac{p_x}{m}\pder{W}{x}
  +\sum_{\ell=0}^{\infty}
  \frac{(-1)^\ell}{(2\ell+1)!}
  \left(\frac{\hbar}{2}\right)^{2\ell}
  \frac{\dd^{2\ell+1}V}{\dd x^{2\ell+1}}
  \frac{\partial^{2\ell+1}W}{\partial p_x^{2\ell+1}}.
  \label{eq:moyal-potential-series}
\end{equation}
当 $V$ 至多为二次多项式时，$V^{(3)}$ 及更高导数全部为零，精确量子传播便与经典 Liouville 传播一致。准确条件是哈密顿量的 Weyl 符号对全部相空间变量至多二次，而不只是“势能看起来像抛物线”。

把 Moyal 方程写成守恒流形式、检验归一化或推导矩方程时都使用了分部积分。对无界相空间，必须要求 $W$ 及所需阶数的导数在无穷远衰减得足够快，使边界通量为零；对周期相空间，边界项因周期性相消。若数值计算只保留有限相空间盒，则吸收、反射或周期边界会改变方程，不能在未检查边界概率质量的情况下沿用这些推导。

\section{非线性怎样产生高阶导数}

考虑四次势
\begin{equation}
  V(x)=\frac{\lambda}{4}x^4.
\end{equation}
因为 $V'''(x)=6\lambda x$ 且五阶以上导数为零，精确方程为
\begin{equation}
  \pder{W}{t}
  =-\frac{p_x}{m}\pder{W}{x}
  +\lambda x^3\pder{W}{p_x}
  -\frac{\hbar^2\lambda x}{4}\frac{\partial^3W}{\partial p_x^3}.
  \label{eq:quartic-wigner-evolution}
\end{equation}
前两项沿经典非线性轨迹剪切分布，第三项不能表示成普通确定性轨迹的概率守恒流。舍弃第三项便得到这一模型的截断 Wigner 演化。

\begin{intuitionbox}
非线性先把一个平滑相空间云团拉伸和折叠，再使分布出现越来越细的结构。精确量子项对这些结构敏感。TWA 在结构仍比量子相空间尺度粗时较有希望，在干涉条纹和回复形成后通常失效。
\end{intuitionbox}

\section{Bopp 算符：另一条推导路线}

与其每次显式展开星积，也可以用 Bopp 对应规则。对坐标变量，
\begin{align}
  \hat x\rhohat&\leftrightarrow
  \left(x+\frac{\ii\hbar}{2}\partial_{p_x}\right)W,
  &\rhohat\hat x&\leftrightarrow
  \left(x-\frac{\ii\hbar}{2}\partial_{p_x}\right)W,\\
  \hat p_x\rhohat&\leftrightarrow
  \left(p_x-\frac{\ii\hbar}{2}\partial_x\right)W,
  &\rhohat\hat p_x&\leftrightarrow
  \left(p_x+\frac{\ii\hbar}{2}\partial_x\right)W.
  \label{eq:bopp-real}
\end{align}
对玻色复振幅，令 $\alpha$ 与 $\alpha^*$ 在微分时相互独立，则
\begin{align}
  \ha\rhohat&\leftrightarrow
  \left(\alpha+\frac12\partial_{\alpha^*}\right)W,
  &\rhohat\ha&\leftrightarrow
  \left(\alpha-\frac12\partial_{\alpha^*}\right)W,\\
  \had\rhohat&\leftrightarrow
  \left(\alpha^*-\frac12\partial_\alpha\right)W,
  &\rhohat\had&\leftrightarrow
  \left(\alpha^*+\frac12\partial_\alpha\right)W.
  \label{eq:bopp-complex}
\end{align}
把哈密顿量或 Lindblad 主方程左右两侧的算符依次替换为\cref{eq:bopp-real,eq:bopp-complex}，即可得到 Wigner 偏微分方程。

\section{从偏微分方程到随机过程}

标准 Fokker--Planck 方程只含一阶和二阶导数：
\begin{equation}
  \pder{W}{t}
  =-\partial_i(A_iW)
  +\frac12\partial_i\partial_j(D_{ij}W).
  \label{eq:fokker-planck}
\end{equation}
当扩散矩阵可以适当分解时，\cref{eq:fokker-planck} 与随机微分方程对应。一阶项给出漂移，二阶项给出演化过程中持续注入的噪声。闭合、幺正、含四次相互作用的玻色系统在 Wigner 表象中通常出现一阶和三阶导数；TWA 丢弃三阶项后只剩确定性漂移，因此量子噪声只通过初始抽样进入。开放系统则可能保留二阶扩散，演化过程中仍有随机噪声。

\begin{warningbox}
“TWA 轨迹是确定性的”只适用于截断后的闭合幺正模型。含 Lindblad 损耗、测量或热浴时，Wigner 方程可以产生真正的二阶扩散，轨迹方程也会含随时间变化的噪声。
\end{warningbox}

\section{矩方程为何有时比整个分布更可靠}

在上一节说明的零边界通量条件下，高阶导数对低阶多项式矩的直接贡献可能因分部积分而消失。例如\cref{eq:quartic-wigner-evolution} 中的三阶 $p_x$ 导数不直接改变 $p_x^0,p_x^1,p_x^2$ 矩。然而，非线性会使低阶矩与更高阶矩耦合，高阶结构随后仍可能反馈回来。因此，某个低阶观测量在一段时间内准确，不代表完整 Wigner 函数或所有高阶关联都同样准确。

\begin{chaptercheck}
看到“经典轨迹”时，检查它来自精确的二次哈密顿量，还是来自对高阶 Moyal 导数的截断；看到“随机轨迹”时，检查噪声只在初态出现，还是由二阶扩散在演化中持续产生。
\end{chaptercheck}

\section*{本章小结}

星积保存了算符乘法和不对易性，Moyal 括号把 von Neumann 方程精确转写为相空间演化。Poisson 括号是其领先项，非线性哈密顿量产生高阶导数。TWA 的核心数学操作正是舍弃无法由普通轨迹表示的高阶项，而 Bopp 算符提供了对具体玻色哈密顿量进行这一推导的实用途径。

\begin{exercises}
  \item 使用\cref{eq:star-product} 计算 $x^2\star p_x^2$，保留全部非零项。
  \item 证明 Moyal 括号满足反对称性，并验证 $\{x,p_x\}_{\mathrm{M}}=1$。
  \item 从\cref{eq:moyal-potential-series} 推导四次势的\cref{eq:quartic-wigner-evolution}。
  \item 使用 Bopp 规则重新推导 $(\had\ha)_W=|\alpha|^2-1/2$。
  \item 对六次势 $V(x)=gx^6/6$，写出 Wigner 演化中最高阶的 $p_x$ 导数。
\end{exercises}
